% Inspired by: 
% 	- https://rahil-makadia.github.io/
%	- https://github.com/jitinnair1/autocv
%----------------------------------------------------------------------------------------
%	DOCUMENT DEFINITION
%----------------------------------------------------------------------------------------

% article class because we want to fully customize the page and not use a cv template
\documentclass[letterpaper,12pt]{article}

%----------------------------------------------------------------------------------------
%	FONT
%----------------------------------------------------------------------------------------

% % fontspec allows you to use TTF/OTF fonts directly
% \usepackage{fontspec}
% \defaultfontfeatures{Ligatures=TeX}

% % modified for ShareLaTeX use
% \setmainfont[
% SmallCapsFont = Fontin-SmallCaps.otf,
% BoldFont = Fontin-Bold.otf,
% ItalicFont = Fontin-Italic.otf
% ]
% {Fontin.otf}

%----------------------------------------------------------------------------------------
%	PACKAGES
%----------------------------------------------------------------------------------------
\usepackage{url}
\usepackage{parskip} 	

%other packages for formatting
\RequirePackage{color}
\RequirePackage{graphicx}
\usepackage[usenames,dvipsnames]{xcolor}
\usepackage[scale=0.9,headsep=0pt]{geometry}
\usepackage{siunitx}
\usepackage{lastpage}
\usepackage{fancyhdr}
\fancypagestyle{firstpage}{
	\fancyhf{}
	\renewcommand{\headrulewidth}{0pt}
	\cfoot{Page {\thepage} of \pageref{LastPage}}
}
\fancypagestyle{regular}{
	\pagestyle{firstpage}
	\rhead{Linyi (Tiger) Hou}
}
\fancypagestyle{lastpage}{
	\pagestyle{regular}
	\rfoot{\footnotesize Last updated: \today}
}
\pagestyle{regular}

%tabularx environment
\usepackage{tabularx}
\newcolumntype{L}[1]{>{\raggedright\arraybackslash}p{#1}}
\newcolumntype{R}[1]{>{\raggedleft\arraybackslash}p{#1}}
% centered version of 'X' col. type
\newcolumntype{H}{>{\centering\arraybackslash}X}
\usepackage{makecell}

%for lists within experience section
% lots of options here for bullet style: https://latex-tutorial.com/bullet-styles/
\usepackage{amssymb} % $\blacktriangleright$
\usepackage{pifont}  % \ding{}
\usepackage{enumitem}
\setlist[itemize]{
	nosep,
	label=\ding{118},%\raisebox{0.5ex}{\tiny$\blacktriangleright$},
	after=\strut, 
	leftmargin=1.5em, 
	itemsep=3pt
}

% experience boilerplates
\NewDocumentEnvironment{WorkExperience}{mmmm +b}{
	\begin{tabularx}{\linewidth}{ @{} L{0.86\linewidth} @{}X@{} r @{} }
		\textbf{#1} & &
		\multirow{2}{*}{\makecell[r]{#2 \\ -- #3}} \\[3.75pt]
		\textit{#4} & & \\[3.75pt]
		\begin{minipage}[t]{\linewidth}
			\begin{itemize}
				#5
			\end{itemize}
		\end{minipage} & &
	\end{tabularx}
}{}

\NewDocumentEnvironment{ResearchExperience}{mmm +b}{
  \begin{tabularx}{\linewidth}{@{} L{0.86\linewidth} @{}X@{} r @{}}
    \textit{#1} & &
    \multirow{2}{*}{\makecell[r]{#2 \\ -- #3}} \\[3.75pt]
    \begin{minipage}[t]{\linewidth}
      \begin{itemize}
        #4
      \end{itemize}
    \end{minipage} & &
  \end{tabularx}
}{}

\NewDocumentEnvironment{ProjectExperience}{mm +b}{
  \begin{tabularx}{\linewidth}{@{} L{0.86\linewidth} @{}X@{} r @{}}
    \textit{#1} & &
    \multirow{2}{*}{\makecell[r]{#2}} \\[3.75pt]
    \begin{minipage}[t]{\linewidth}
      \begin{itemize}
        #3
      \end{itemize}
    \end{minipage} & &
  \end{tabularx}
}{}

%custom \section
\usepackage{titlesec}				
\usepackage{multicol}
\usepackage{multirow}

%CV Sections inspired by: 
%http://stefano.italians.nl/archives/26
\titleformat{\section}{\Large\scshape\raggedright}{}{0em}{}[\titlerule]
\titlespacing{\section}{0pt}{10pt}{10pt}
\titleformat{\subsection}{\large\scshape\raggedright}{}{0em}{}
\titlespacing{\subsection}{0pt}{5pt}{5pt}

%for publications
\usepackage[backend=biber,style=ieee,sorting=none,style=numeric,maxbibnames=99]{biblatex}
\setlength{\labelnumberwidth}{24pt}

% no numbering
% \DeclareFieldFormat{labelnumberwidth}{}
% \setlength{\biblabelsep}{0pt}

% reverse numbering (https://tex.stackexchange.com/a/22770)
	% Count total number of entries in each refsection
\AtDataInput{%
  \csnumgdef{entrycount:\therefsection}{%
    \csuse{entrycount:\therefsection}+1}}
	% Print the labelnumber as the total number of entries in the
	% current refsection, minus the actual labelnumber, plus one
\DeclareFieldFormat{labelnumber}{\mkbibdesc{#1}}
\newrobustcmd*{\mkbibdesc}[1]{%
  \number\numexpr\csuse{entrycount:\therefsection}+1-#1\relax}

%Setup hyperref package, and colours for links
\usepackage[unicode, draft=false]{hyperref}
\definecolor{linkcolour}{rgb}{0,0.2,0.6}
\hypersetup{colorlinks,breaklinks,urlcolor=linkcolour,linkcolor=linkcolour}
\addbibresource{pub_journal.bib}
\addbibresource{pub_conference.bib}
\setlength\bibitemsep{5pt}

%for social icons
\usepackage{fontawesome5}

%flag for digital CV
\newif\ifdigital

%----------------------------------------------------------------------------------------
%	BEGIN DOCUMENT
%----------------------------------------------------------------------------------------
\begin{document}
	
	% non-numbered pages
	\pagestyle{empty} 
	


	%----------------------------------------------------------------------------------------
	%	TITLE
	%----------------------------------------------------------------------------------------
	\begin{tabularx}{\linewidth}{@{} H @{}}
		\Huge{Linyi (Tiger) Hou} \\[7.5pt]
		\href{mailto:linyih2@illinois.edu}{\raisebox{-0.05\height}\faEnvelope \ linyih2@illinois.edu}
		\ $|$ \ \href{https://github.com/TigerHou2}{\raisebox{-0.05\height}\faGithub \ github.com/TigerHou2}
		\ $|$ \ U.S. Permanent Resident \\
	\end{tabularx}



	%----------------------------------------------------------------------------------------
	%	EDUCATION
	%----------------------------------------------------------------------------------------
	\section{Education}
	\begin{tabularx}{\linewidth}{@{}l X@{}}
		\textbf{Ph.D. in Aerospace Engineering} & \hfill May 2025 \\
		\multicolumn{2}{@{}l@{}}{University of Illinois Urbana-Champaign \hfill GPA: 3.94/4.00} \\
		\multicolumn{2}{@{}l@{}}{Thesis: Autonomous Lost in Space and Time State Determination using Periodic Variable Stars} \\[3.75pt]
		\textbf{B.S. in Aerospace Engineering, Minor in Computer Engineering} & \hfill May 2021 \\
		\multicolumn{2}{@{}l@{}}{University of Illinois Urbana-Champaign \hfill GPA: 3.93/4.00} \\
	\end{tabularx}
	
	

	%----------------------------------------------------------------------------------------
	%	SKILLS
	%----------------------------------------------------------------------------------------
	\section{Skills}
	\begin{tabularx}{\linewidth}{@{}l X@{}}
		Technical  &  batch filtering, Kalman filtering, Monte Carlo simulation, celestial navigation, orbit determination, spacecraft navigation modeling \\[3.75pt]
		Collaboration  &  technical/proposal writing, MS Office suite, \LaTeX, Git, GitHub, Trello \\[3.75pt]
		Programming  &  Python, C/C++, MATLAB/Simulink, GMAT, Javascript, Julia \\[3.75pt]
	\end{tabularx}


	
	%----------------------------------------------------------------------------------------
	% RESEARCH
	%----------------------------------------------------------------------------------------
	
	\section{Research}
	
	\textbf{Astrodynamics and Planetary Exploration Group (APEX)} 
	\hfill Advisor: Dr. Siegfried Eggl
	
	\begin{ResearchExperience}
		{Interstellar Navigation and Star Identification}
		{Sep 2024}{May 2025}
			\item Developed a star identification technique using C++ which mitigates stellar aberration effects at relativistic speeds, improving star identification success rate from 38\% to 91\%.
	\end{ResearchExperience}

	\begin{ResearchExperience}
		{Celestial Navigation to Supplement DSN Tracking}
		{Jul 2024}{Mar 2025}
			\item Collaborated with colleagues to create navigation simulation tools in Python to assess the feasibility of reducing Deep Space Network tracking cadence by using onboard celestial navigation for the OSIRIS-APEX and CAPSTONE missions.
			\item Designed sensor scheduling algorithms for celestial navigation which reduced navigation uncertainty by over 10\% compared to existing strategies in Monte Carlo simulation. 
			\item Implemented spherical harmonics for an open source orbit propagator in C++.
	\end{ResearchExperience}

	\begin{ResearchExperience}
		{Lost in Space and Time Navigation using Variable Stars}
		{Nov 2022}{Mar 2025}
			\item Developed an autonomous optical navigation concept using variable star light curves which delivers timing accuracy within 1 second without needing prior position/time information.
			\item Simulated star tracker observations and demonstrated the feasibility of reestablishing Earth pointing via the proposed method through Monte Carlo analysis.
	\end{ResearchExperience}
	

	\textbf{Aerospace Mission Analysis Laboratory} 
	\hfill Advisor: Dr. Zachary R. Putnam
	
	\begin{ResearchExperience}
		{Lost in Space X-ray Pulsar Navigation (XNAV)}
		{Jun 2021}{Dec 2023}
			\item Developed algorithms to resolve spacecraft position and time from pulsar timing data.
			\item Created a photon time of arrival simulation tool in Python and validated simulation results against the pulsar timing package PINT to within \qty{50}{\micro\s}.
	\end{ResearchExperience}

	\begin{ResearchExperience}
		{Velocity-based Initial Orbit Determination}
		{Aug 2019}{Apr 2024}
			\item Developed a novel initial orbit determination technique using batched sequential range-rate measurements that achieved over 10x better accuracy than state-of-the-art methods.\hspace{-1em}
	\end{ResearchExperience}


	\textbf{Space Systems Optimization Laboratory} 
	\hfill Advisor: Dr. Koki Ho
	
	\begin{ResearchExperience}
		{Architecture Trade Studies on ISRU Technologies for Human Space Exploration}
		{Sep 2018}{Sep 2019}
			\item Worked with two graduate students to discuss the development of a mixed-integer linear programming framework for space logistics optimization.
			\item Conducted literature review of in-situ resource utilization (ISRU) technologies to quantify input/output rates of technology options for an architecture optimization tool.
	\end{ResearchExperience}
	

	
	%----------------------------------------------------------------------------------------
	% WORK
	%----------------------------------------------------------------------------------------
	
	\section{Work Experience}
	
	\begin{WorkExperience}
		{Teaching Assistant}
		{Jun 2021}{May 2024}
		{University of Illinois Urbana-Champaign (AE 202, AE 443, AE 483)}
			\item Worked with 2-5 faculty and staff to teach lessons, hold office hours, and improve curriculum for courses in flight mechanics, systems engineering, and UAV controls.
			\item Helped 20+ students formulate and implement UAV state-space model controllers with various sensors, e.g., laser rangefinders, optical flow cameras, and time-of-flight anchors.
			\item Assisted in developing LQR control-related code and instructional materials for a UAV course using C and Python.
	\end{WorkExperience}

	\begin{WorkExperience}
		{Software Engineering Intern}
		{May 2019}{Aug 2019}
		{Collins Aerospace}
			\item Migrated an avionics data interpreter from 16-bit to 32-bit using C++ and Python.
			\item Automated weekly builds and integration in TortoiseSVN using Python.
	\end{WorkExperience}
	

	
	%----------------------------------------------------------------------------------------
	% PROJECTS
	%----------------------------------------------------------------------------------------

	\section{Projects}

	\begin{ProjectExperience}
		{Collaborative Projects in Flight Dynamics}
		{2022, 2023\\\strut}
			\item Planetary entry: Created a 3-DOF entry simulation to assess the flight performance of several novel Apollo-derived entry guidance algorithms at Mars.
			\item Trajectory optimization: Applied optimal control theory to identify time- and fuel-optimal low-thrust Earth-Mars transfer trajectories. 
	\end{ProjectExperience}

	\begin{ProjectExperience}
		{Team Lead/Advisor, NASA's RASC-AL Competition}
		{2019, \\ 2022, 2023}
			\item Led a team of 10+ students to develop a reusable crewed lunar lander concept (structure, power, thermal, life support, orbital mechanics, budget, etc.). 
			\item Advised two student teams who developed technical requirements and performed trade studies for lunar and Martian in-situ resource utilization architectures. 
	\end{ProjectExperience}



	%----------------------------------------------------------------------------------------
	%	AWARDS/AFFILIATIONS
	%----------------------------------------------------------------------------------------
	\section{Honors and Awards}

	{\setlength{\parskip}{0.25\baselineskip}

	\begin{tabularx}{\linewidth}{@{}l X@{}}
		\textbf{Aerospace Engineering Alumni Advisory Board Fellowship} & \hfill 2024 \\
		Department of Aerospace Engineering, University of Illinois Urbana-Champaign & \hfill \\
	\end{tabularx}

	\begin{tabularx}{\linewidth}{@{}l X@{}}
	\textbf{Best Student Paper, 46\textsuperscript{th} AAS GN\&C Conference} & \hfill 2024 \\
	American Astronautical Society (AAS) Rocky Mountain Section & \hfill \\
	\end{tabularx}
	
	\begin{tabularx}{\linewidth}{@{}l X@{}}
		\textbf{Yee Fellowship} & \hfill 2021, 2024 \\
		Grainger College of Engineering, University of Illinois Urbana-Champaign & \hfill \\
	\end{tabularx}

	\begin{tabularx}{\linewidth}{@{}l X@{}}
	\textbf{David and Catherine Thompson Space Technology Scholarship} & \hfill 2020 \\
	American Institute of Aeronautics and Astronautics (AIAA) & \hfill \\
	\end{tabularx}

	\begin{tabularx}{\linewidth}{@{}l X@{}}
	\textbf{H.S. Stillwell Problem Solving Award} & \hfill 2020 \\
	Department of Aerospace Engineering, University of Illinois Urbana-Champaign & \hfill \\
	\end{tabularx}

	\begin{tabularx}{\linewidth}{@{}l X@{}}
	\textbf{Robert W. McCloy Memorial Award} & \hfill 2019 \\
	Department of Aerospace Engineering, University of Illinois Urbana-Champaign & \hfill \\
	\end{tabularx}

	\begin{tabularx}{\linewidth}{@{}l X@{}}
	\textbf{James Scholar} & \hfill 2018 - 2021 \\
	University of Illinois Urbana-Champaign & \hfill \\
	\end{tabularx}

	\begin{tabularx}{\linewidth}{@{}l X@{}}
	\textbf{Dean's List} & \hfill 2017 - 2021 \\
	University of Illinois Urbana-Champaign & \hfill \\
	\end{tabularx}
	
	}



	%----------------------------------------------------------------------------------------
	%	PUBLICATIONS
	%----------------------------------------------------------------------------------------

	\begin{refsection}[pub_journal.bib]
		\DeclareFieldFormat{labelnumber}{J.\mkbibdesc{#1}}
		\nocite{*}
		\printbibliography[title={Journal Articles}]
	\end{refsection}
	
	\begin{refsection}[pub_conference.bib]
		\DeclareFieldFormat{labelnumber}{C.\mkbibdesc{#1}} 
		\nocite{*}
		\printbibliography[title={Conference Proceedings}]
	\end{refsection}
	
	%----------------------------------------------------------------------------------------
	%	LAST UPDATED
	%----------------------------------------------------------------------------------------
	\vfill
	\center{\footnotesize Last updated: \today}
	
\end{document}